\def\stat{filipov}





\def\tit{МЕТОД ПОВЫШЕНИЯ ПЕРТИНЕНТНОСТИ ИНФОРМАЦИИ 


В~РЕКОМЕНДАТЕЛЬНЫХ СИСТЕМАХ ПОДДЕРЖКИ 


ЖИЗНЕОБЕСПЕЧЕНИЯ НА~ОСНОВЕ НЕЯВНЫХ ДАННЫХ$^*$}





\def\titkol{Метод повышения пертинентности информации 


в %рекомендательных 


системах поддержки 


жизнеобеспечения} % на~основе неявных данных}








\def\aut{С.\,А.~Филиппов$^1$, В.\,Н.~Захаров$^2$}





\def\autkol{С.\,А.~Филиппов, В.\,Н.~Захаров}





\titel{\tit}{\aut}{\autkol}{\titkol}





\index{Филиппов С.\,А.}


\index{Захаров В.\,Н.}


\index{Philippov S.\,A.}


\index{Zakharov V.\,N.}





{\renewcommand{\thefootnote}{\fnsymbol{footnote}}


\footnotetext[1]


{Работа выполнена при поддержке Министерства образования и~науки РФ (уникальный идентификатор 


проекта RFMEFI60414X0139).}}





\renewcommand{\thefootnote}{\arabic{footnote}}


\footnotetext[1]{Институт проблем информатики Федерального исследовательского


  центра <<Информатика и~управление>> Российской академии наук, 


\mbox{stanislav@philippov.ru}}


\footnotetext[2]{Институт проблем информатики Федерального исследовательского


  центра <<Информатика и~управление>> Российской академии наук, 


\mbox{vzakharov@ipiran.ru}}








\label{st\stat}





                  \thispagestyle{headings}


                  


                  \vspace*{-18pt}





  


   


  \Abst{Сегодня, когда методы релевантного поиска практически достигли своих пределов, 


все большее внимание уделяется повышению пертинент\-ности информации. Особенно это 


справедливо в~области систем поддержки жизнеобеспечения, связанных с~розничной 


торговлей товарами и~(или) услугами в~Интернете, где происходит серьезная борьба за 


интерес покупателя. Практически все современные крупные ин\-тер\-нет-ма\-га\-зи\-ны стараются 


адаптировать внешний вид и~содержимое своих страниц под нужды конкретных 


пользователей, в~основном с~применением коллаборативной фильтрации (collaborative 


filtering, CF). Преимуществами данного подхода являются увеличение вовлеченности 


посетителей, улучшение пользовательского опыта и~повышение конверсии.


  В~данной работе представлено описание метода повышения пертинентности 


информации, опирающегося на неявные данные, т.\,е.\ на результаты обработки активности 


пользователей, связанных с~процессом принятия решения, а~не с~его результатом (как 


в~случае обработки явных данных). В~основе метода лежит комбинированное 


использование подходов Item--Item CF и~User--User CF, позволяющее предлагать посетителю 


пертинентное информационное предложение даже в~ситуациях, когда сведения 


о~пользовательской активности отсутствуют или малоинформативны. Возможности 


предложенного метода проверены с~помощью экспериментального образца программного 


комплекса повышения пертинентности информации, установленного на действующий 


интернет-магазин Thaisoap. Полученные результаты подтверждают, что даже ограниченное 


применение метода достаточно серьезно для электронной коммерции улучшает 


пользовательский опыт и~доходы владельца.}


   


  \KW{пертинентность; коллаборативная фильтрация; интернет-магазин; рекомендательная 


система; неявные данные}





\DOI{10.14357\!/\!08696527160401} 








\vskip 8pt plus 9pt minus 6pt








      \thispagestyle{headings}


      


                        \vspace*{-16pt}


  


\section{Введение}





\vspace*{-2pt}





  Одним из современных трендов в~развитии сети Интернет является 


персонализация~--- способ подачи информации на сайтах, который 


подстраивается под уникальные потребности конкретных посетителей. 


Современные поисковые системы, социальные сети, форумы, новостные 


ресурсы и~системы поддержки жизнеобеспечения (здесь и~далее под такими 


системами подразумеваются системы электронной коммерции,  


ин\-тер\-нет-ма\-га\-зи\-ны) стараются адаптировать внешний вид и~содержимое 


(контент) своих страниц под нужды конкретных пользователей. По результатам 


исследования компании Evergage, в~2015~г.\ персонализацию в~реальном 


времени использовали 44\% веб-сай\-тов, 17\% мобильных сайтов,  


13\%~веб-при\-ло\-же\-ний и~9\%~мобильных приложений~[1]. При 


этом~78\%~тех, кто не использует персонализацию сейчас, утверждают, что 


планируют начать ее использовать в~течение следующих 12~мес. Увеличение 


вовлеченности посетителей, улучшение пользовательского опыта и~повышение 


конверсии считаются самыми важными результатами ее применения. 


  


  В сфере электронной коммерции основным инструментом персонализации 


контента служат рекомендательные системы, обеспечивающие автоматическую 


обработку данных о пользовательской активности и~выработку рекомендаций 


на товары и~услуги, которые могут быть интересны конкретным 


пользователям~[2]. Исходными данными для анализа поведения пользователей 


являются сведения об их активности, которые могут собираться явным или 


неявным образом. Явным образом получают персональные данные при 


регистрации, покупках, вы\-став\-ле\-нии пользователем оценок тем или иным 


объектам на сайтах, в~результате голосований и~опросов, другие данные, 


связанные с~уже принятыми пользователем решениями. 


  %


  При этом сегодня практически не уделяется внимания неявным данным, 


таким как, например, время пребывания посетителей на отдельных страницах, 


копирование данных со страниц, выделение отдельных характеристик товаров 


и~т.\,п. Зачастую это связано с~тем, что такие данные составляют огромные 


массивы, которые являются неоднородными и~достаточно сложными для 


интерпретации. Более того, их хранение в~реляционных базах данных 


фактически невозможно~[3]. Поэтому, как правило, используются так 


называемые NoSQL системы управления данными (HBase, Cassandra). Их 


характерными особенностями являются отказ от транзакций, практически 


линейная масштабируемость, высокая скорость обработки запросов, отсутствие 


жесткой схемы данных. 


  


  Необходимо отметить, что адаптация под конкретного пользователя~--- 


весьма сложная задача, поскольку для ее решения необходимо принимать во 


внимание как присущие человеку неопределенность и~спонтанность в~рамках 


конкретного ин\-тер\-нет-ре\-сур\-са, так и~множество неопределенностей, 


связанных с~особенностями функционирования сети Интернет. В~контексте 


проблемы персонализации контента (а~также прогнозирования, выявления 


предпочтений и~групп схожих ресурсов) встает задача обработки собранных 


данных и~выявления определенных закономерностей, позволяющих сделать 


выводы о конкретных предпочтениях пользователей. Таким образом, основной 


целью обработки данных о пользовательской активности является извлечение 


полезной информации, которая может, в~свою очередь, использоваться для 


решения следующих задач~[4]:


  \begin{enumerate}[(1)]


\item кластеризация ресурсов. Группирование схожих по множеству 


посетителей ресурсов в~несколько кластеров (групп) ресурсов. 


Кластеризация позволяет строить каталоги ресурсов, а~также выявлять 


недостатки существующих тематических каталогов;


\item кластеризация пользователей. Группирование схожих пользователей 


в~клас\-те\-ры аналогично кластеризации ресурсов. Позволяет выявлять 


группы пользователей со схожими интересами;


\item построение устойчивых поведенческих профилей пользователей в~виде 


перечня групп ресурсов, посещаемых как данным пользователем, так 


и~схожими с~ним пользователями; 


\item построение расширенных профилей пользователей, включающих  


со\-ци\-аль\-но-де\-мо\-гра\-фи\-че\-ские данные (анкеты), описательные 


статистики и~поведенческие профили. Расширенные профили позволяют 


классифицировать новых пользователей, выявлять зависимости между 


пользовательским поведением и~со\-ци\-аль\-но-де\-мо\-гра\-фи\-че\-ски\-ми 


характеристиками; 


\item сегментация клиентской базы на основе расширенных профилей 


позволяет выделять сегменты как по анкетным данным клиентов, так и~по их 


поведению. Эта информация используется при маркетинговых 


исследованиях;


\item персонализация контента. Представление каждому пользователю сайта 


наиболее интересной для него информации в~наиболее удобном для него 


виде. Знание информационных предпочтений пользователя позволяет 


динамически перестраивать контент сайта;


\item построение карт сходства ресурсов и~пользователей. Позволяет 


отображать множества наиболее посещаемых ресурсов и~наиболее активных 


пользователей в~виде точечного графика. Схожим ресурсам (пользователям) 


соответствуют близкие точки на карте. Карту сходства можно использовать 


как графическое средство навигации.


\end{enumerate}





  Существуют различные методы и~подходы, используемые на практике при 


решении перечисленных выше задач. Весь класс этих методов, которые 


принято называть методами коллаборативной фильтрации, и~рассмотрен далее.





\vspace*{-6pt}


   


\section{Методы коллаборативной фильтрации}





\vspace*{-2pt}





  В основе методов коллаборативной фильтрации 


лежит предположение о консервативности пользовательских предпочтений 


(т.\,е.\ пользователи, одинаково оценивающие определенные объекты, скорее 


всего, аналогичным образом будут оценивать и~новые объекты со сходными 


характеристиками)~[5]. По существу, рекомендации базируются на 


автоматическом сотрудничестве множества пользователей и~на выделении 


(методом фильтрации) тех пользователей, которые демонстрируют схожие 


предпочтения или шаблоны поведения. Таким образом, методы 


коллаборативной фильтрации вырабатывают рекомендации, основанные на 


модели предшествующего поведения пользователя и~с учетом поведения 


пользователей со схожими характеристиками. 


  


  Наибольшее распространение в~сфере электронной коммерции получили 


рекомендательные системы, использующие следующие методы 


коллаборативной фильтрации: 


  \begin{itemize}


\item коллаборативная фильтрация посредством анализа предпочтений групп 


пользователей со схожими интересами (User--User CF);


\item коллаборативная фильтрация посредством анализа взаимосвязей между 


информационными единицами (Item--Item CF).


  \end{itemize}


   При этом по отдельности методы Item--Item CF и~User--User CF используют 


только часть собираемой рекомендательными системами информации 


о~пользовательской активности и~предпочтениях при выработке рекомендаций. 


Более эффективным может быть использование гибридных подходов. Так, 


например, в~работе~[6] предлагается алгоритм комбинированной фильтрации, 


основная идея которого заключается в~получении оценки неизвестного 


рейтинга как взвешенной суммы оценок на основании фильтрации по 


транзакциям, фильтрации по товарам и~смешанной фильтрации (на основании 


рейтингов похожих товаров в~похожих транзакциях).





\vspace*{-6pt}


  


\section{Обзор метода повышения пертинентности информации}





\vspace*{-2pt}





  Авторами данной работы по итогам проведенных прикладных научных 


исследований предлагается собственный метод повышения пертинентности 


информации, сочетающий сильные стороны обоих рассмотренных выше 


методов коллаборативной фильтрации и~позволяющий более полно 


использовать доступную информацию о пользователях и~информационных 


единицах, в~первую очередь неявные данные о пользователе. 


  


  Суть предлагаемого метода заключается в~совместном использовании при 


выработке рекомендаций оценок, полученных с~использованием как метода 


Item--Item CF, так и~метода User--User CF. Для новых пользователей или 


пользователей с~нерепрезентативной историей посещений предлагается 


генерировать рекомендации, базируясь на данных о подобии информационных 


единиц (метод Item--Item CF). Таким образом решается проблема <<холодного 


старта>> и~повышается качество рекомендаций для малоактивных 


пользователей, а~также пользователей со слабо выраженными предпочтениями. 


По мере накопления данных о пользовательских предпочтениях 


и~формировании его поведенческого профиля приоритет предлагается отдавать 


оценкам, полученным с~использованием метода User--User CF, исходя из 


гипотезы о том, что при наличии качественных поведенческих профилей 


пользователей данный метод позволяет более точно предсказывать их 


предпочтения. При этом в~случае необходимости рекомендации, полученные 


с~использованием метода User--User CF, могут дополняться предложениями 


информационных единиц, полученными на основе их рейтинга популярности.


  


  Высказанные гипотезы и~работоспособность метода были экспериментально 


исследованы в~составе действующего ин\-тер\-нет-ма\-га\-зи\-на Thaisoap. 


Магазин ориентирован на продажу натуральной тайской косметики и~


кокосового масла. Каталог товаров магазина содержит 


более~1\,500~наименований товаров, разбитых на~180~классов (44~корневых 


класса, 136~подклассов). Ежедневно по данным за~2015~г.\ магазин посещают 


в~среднем около~1\,500 посетителей и~проводят на нем в~среднем~11~мин 


каждый (на каждого посетителя приходится в~среднем~28~переходов по 


ссылкам). Для теоретических исследований использовалась тестовая выборка 


данных за~IV~квартал 2015~г., экспериментальные исследования проведены 


в~фев\-ра\-ле--мае 2016~г. В~обоих случаях каталог товаров был неизменен. 


Дополнительно использовались статистические данные аналитической системы 


Яндекс.Метрика ({\sf http://metrika.yandex.ru}).


  


\vspace*{-6pt}





\section{Определение подобия информационных единиц 


по~неявным пользовательским предпочтениям}





\vspace*{-2pt}





  В целях решения задачи формирования рекомендации с~уместной 


информацией в~условиях недостаточности знаний о пристрастиях 


пользователей авторами предлагается использовать метод, в~основе которого 


лежит расчет близости пар и~последующая группировка (кластеризация) 


информационных единиц на основе данных пользователей, последовательно 


просматривающих несколько товаров. При отсутствии данных предлагается 


использовать обычные классификаторы с~учетом цены и~параметров объектов, 


список <<Новинки>>, а~также матрицу <<С~этим товаром покупают>> 


(аксессуары, дополняющие основную покупку). При данном подходе явное 


участие пользователей ин\-тер\-нет-ма\-га\-зи\-на в~формировании рейтинга 


товаров не требуется.


  


  Первым шагом алгоритма является построение матрицы подобия 


информационных единиц, где и~по вертикали, и~по горизонтали присутствуют 


все информационные единицы интернет-магазина. Заполнение матрицы 


происходит по следующему правилу: если пользователь последовательно 


просмотрел два товара, то вес подобия в~матрице для этих двух товаров 


увеличивается на~1. Для обработки матрицы в~целях выявления групп 


информационных единиц, которые являются близкими по своим оценкам 


подобия, из всех известных алгоритмов кластеризации в~результате 


проведенного моделирования был выбран современный производительный 


алгоритм Affinity Propagation. Одним из преимуществ данного алгоритма 


является отсутствие необходимости предварительной оценки оптимального 


количества кластеров~[7].


  


  С помощью статистического пакета R была проведена обработка тестового 


массива данных ин\-тер\-нет-ма\-га\-зи\-на Thaisoap: построена матрица 


подобия по всему доступному временн$\acute{\mbox{о}}$му периоду и~проведена ее обработка 


по алгоритму Affinity Propagation.


  


  Всего алгоритм выделил~64~кластера, наиболее крупными из которых 


оказались кластеры с~номерами~5 (75~объектов), 8~(44~объекта), 


10~(30~объектов), 19~(27~объектов) и~55~(31~объект). Качество работы 


алгоритма можно оценить на примере кластера номер~5, описание которого 


представлено в~таблице. 


%  


  В~частности, видно, что для референсной информационной единицы 


<<массажное кокосовое масло>> в~кластер подобия попали товары на основе 


кокосового масла или косвенно ассоциирующиеся с~кремами и~маслами для 


ухода за телом.


  


  \begin{table}\small


  \begin{center}


  \begin{tabular}{|p{12mm}|p{24mm}|p{80mm}|}


  \multicolumn{3}{c}{Детализация кластера номер 5}\\


  \multicolumn{3}{c}{\ }\\[-6pt]


  \hline


\multicolumn{1}{|c|}{Кластер}&\multicolumn{1}{c|}{


\tabcolsep=0pt\begin{tabular}{c}Референсная\\ информационная\\ 


единица\end{tabular}}&\multicolumn{1}{c|}{Примеры товаров из кластера}\\


\hline


ID: 5\newline


Size: 75&ID: 76.\newline


Не\-ра\-фи\-ни\-ро\-ван\-ное 100\% мас\-саж\-ное ко\-ко\-со\-вое мас\-ло <<Citronella>> Tropicana, 100~мл&ID: 


43.\newline


Кокосовое мас\-ло Tropicana, 1~л, нерафинированное\newline


ID: 51.\newline


Кокосовое масло нерафинированное Tropicana в~аптекарском флаконе, 90~мл\newline


ID: 466.\newline


Восстанавливающий кокосовый ЛОСЬОН для тела Tropicana <<Sweet Coconut>> (без 


парабенов), 200~мл\newline


ID: 624.\newline


Мас\-ка-экс\-фо\-ли\-ант для лица <<Морской кол\-ла\-ген>> Artiscent, 100~мл\newline


ID: 1234.\newline


Мининабор шампунь и~кондиционер для волос <<Золотой шелк с~экстрактом 


шелковицы>>\\


\hline


\end{tabular}


\end{center}


\vspace*{-9pt}


\end{table}


  


  Рассмотрение других полученных в~результате работы алгоритма кластеров 


показывает, что в~отдельные группы были выделены товары, имеющие схожие 


потребительские характеристики. Так, например, кластер номер 55 (<<Синий 


тайский бальзам от варикоза>>) сформирован преимущественно бальзамами 


(<<Традиционный малый КРАСНЫЙ тайский бальзам (для целебного массажа) 


Korn Herb>> и~т.\,п.)\ и~лаками (<<Противогрибковый лак <<Демиктен>> 


и~т.\,п.)\ име\-ющи\-ми ле\-чеб\-но-про\-фи\-лак\-ти\-че\-скую направленность. 


В~кластер номер~10 (<<Набор <<Здоровые волосы>> Tropicana>>) попали 


товары для волос и~тела (<<Сыворотка для волос <<Romance>>>>, <<Набор 


<<Роскошные волосы>> Tropicana>> и~т.\,п.). В~качестве референсной 


информационной единицы может выступать товар, который заинтересовал 


потенциального покупателя в~каталоге ин\-тер\-нет-ма\-га\-зи\-на. Предложения 


же рекомендательной системы в~этом случае будут формироваться, используя 


товары с~оценками подобия, наиболее близкими к~референсному товару (т.\,е.\ 


входящие в~соответствующий кластер подобия). Таким образом, если новый 


посетитель выберет в~каталоге массажное кокосовое масло <<Citronella>> 


Tropicana, то рекомендательная система предложит ему посмотреть кокосовое 


масло Tropicana, кокосовое масло нерафинированное Tropicana в~аптекарском 


флаконе, восстанавливающий кокосовый ЛОСЬОН для тела Tropicana и~т.\,д. 


В~зависимости от конфигурации рекомендательного сервиса предложения 


товаров (на базе метода Item--Item CF) могут дополняться наиболее 


популярными на текущий момент товарами или новыми товарами, 


требующими продвижения (например, из пяти слотов для рекомендованных 


товаров один может быть отведен для продвижения новых продуктов из 


линейки массажных масел). 


  


  В целом необходимо отметить, что использование метода Item--Item CF, 


основанного на оценке близости информационных единиц, позволяет получить 


рекомендации, интуитивно понятные посетителям магазина, не имея данных об 


их предпочтениях. Так, например, при выборе бальзама от варикоза 


рекомендательная система предложит посмотреть традиционный бальзам для 


целебного массажа и~противогрибковый лак, что в~целом будет достаточно 


точно соответствовать заданной посетителями тематике поиска  


(ле\-чеб\-но-про\-фи\-лак\-ти\-че\-ские средства на основе кокосового масла).


  


\vspace*{-6pt}





\section{Формирование поведенческих профилей и~выявление 


групп пользователей со~схожими характеристиками}





\vspace*{-2pt}





  По мере накопления данных о пользовательской активности появляется 


возможность формирования поведенческих профилей и~использования метода 


User--User CF для генерации рекомендаций, более полно соответствующих 


предпочтениям конкретных пользователей. В~результате проведенного 


имитационного моделирования было установлено, что наиболее подходящим 


с~учетом существующих для работы ограничений по скорости и~объемам 


обработки данных для выявления групп пользователей со схожими 


предпочтениями (User--User CF) является метод кластеризации К-сред\-них  


(K-means). В целом данный статистический метод прост в~реализации и~


является хорошо масштабируемым~[8]. Вычислительная сложность алгоритма~--- 


$O(nkl)$, где $n$~--- число объектов; $k$~--- число кластеров; $l$~--- число 


итераций. Одним из недостатков данного метода является необходимость 


заранее задавать число кластеров для разбиения. Тем не менее оптимальное для 


разбиения число кластеров достаточно просто получить посредством 


минимизации суммы внутрикластерных расстояний. Другим не\-до\-стат\-ком 


алгоритма К-сред\-них является прямая зависимость производительности 


алгоритма от числа итераций. Число итераций может быть уменьшено заданием 


близких к~оптимальным начальных значений центроидов кластеров (данный 


подход используется в~модифицированном методе К-сред\-них~--- K-means++).


  


  Первым шагом алгоритма является построение матрицы активности 


пользователей. В~тестовом массиве данных было доступно только количество 


обращений к~конкретной категории товаров. Таким образом, каждая строка 


матрицы активности представляет собой вектор оценок, соответствующих 


различным категориям товаров (тематический профиль пользователя). Профиль 


пользователя характеризует степень его интереса к~каждой группе товаров. 


  


  Следующим шагом (перед выявлением групп пользователей со схожими 


характеристиками методами кластеризации) является обработка полученной 


матри-\linebreak\vspace*{-12pt}





\pagebreak





\





\vspace*{-12pt}


\begin{floatingfigure}{77mm} %fig1


   \vspace*{1pt}


 \begin{center}


 \mbox{%


 \epsfxsize=69.58mm


\epsfbox{fil-1.eps}


 }


\end{center}


 \vspace*{-9pt}


\Caption{Гистограмма расстояний}


%\begin{figure} %fig2


  \vspace*{12pt}


 \begin{center}


 \mbox{%


 \epsfxsize=73.558mm


\epsfbox{fil-2.eps}


 }


\end{center}


 \vspace*{-9pt}


\Caption{Анализ количества кластеров}


%\end{figure}


\vspace*{6pt}


\end{floatingfigure}





\noindent


цы активности с~целью вычисления попарных расстояний между ее 


элементами. В~качестве метрики рас\-сто\-яния (функция сходства) было 


использовано рас\-сто\-яние Евклида (геометрическое рас\-сто\-яние в~многомерном 


пространстве), которое является одной из наиболее простых для реализации 


и~часто используемых на практике метрик на сегодняшний день. На рис.~1 


пред\-став\-ле\-на гистограмма рас\-сто\-яний, полученная в~результате обработки 


матрицы с~попарными рас\-сто\-яни\-ями между объектами с~по\-мощью 


статистического пакета~R. Предварительное рассмотрение результатов 


поз\-во\-ля\-ет сделать вывод о~том, что предпочтения (векторы активности) 


пользователей ин\-тер\-нет-ма\-га\-зи\-на Thaisoap не сильно отличаются (т.\,е., 


скорее всего, большинство пользователей интересуются схожими категориями 


товаров).


   





  Одним из исходных параметров для применения метода К-сред\-них является 


число кластеров, на которое необходимо разбить исследуемый массив данных. 


Для определения оптимального числа кластеров было проведено специальное 


исследование. В~качестве критерия использовалась внутрикластерная сумма 


квадратов расстояний, которую необходимо было минимизировать. На рис.~2 


приведены результаты расчета внутрикластерной суммы квадратов расстояний 


по методу локтя (Elbow method). В~результате исследования тестового набора 


данных было определено оптимальное для данного тестового массива число 


кластеров, равное~30.


  


  





  Таким образом, были получены все необходимые параметры для проведения 


кластеризации данных по матрице активности пользователей. Результат работы 


алгоритма кластеризации позволил выделить несколько наиболее крупных 


клас\-те\-ров пользователей с~номерами~1, 9, 20, 21 и~22. Пользователи из кластера 


под номером~1 (размер кластера~1384) демонстрируют слабое предпочтение ко 


всем категориям товаров. Возможно, что в~кластер с~номером~1 попали 


посетители сайта, которые пришли без конкретной цели, например просто 


ознакомиться с~предлагаемым ассортиментом товаров. Пользователи из 


кластера номер~9 (размер~180) демонстрируют явное предпочтение 


к~категории cat9 (<<Лицо>>). Пользователи из кластера номер~20 


(размер~1366) демонстрируют предпочтение к~категории cat7 (<<Кокосовое 


масло>>). Сумма квадратов расстояний относительно других кластеров мала, 


что говорит о небольшом различии объектов внутри кластера. Пользователи из 


кластера номер~21 (размер~622) демонстрируют также предпочтение 


к~категории cat7 (<<Кокосовое масло>>). Пользователи из кластера номер~22 


(размер~180) демонстрируют предпочтение к~категориям cat47 (<<MUST 


HAVE Зима 2016>>) и~cat49 (<<ХИТЫ нашего магазина>>). Проведение 


расчетов на основании данных за другие недели дало похожие результаты.


  


  Таким образом, данные, полученные в~результате обработки исходных данных 


о~пользовательской активности кластера, могут быть использованы для 


формирования предложений рекомендательной системой. Качество 


рекомендаций будет зависеть от адекватности данных о пользовательских 


предпочтениях (другими словами, качества сформированных поведенческих 


профилей). Рекомендоваться пользователям будут информационные единицы, 


имеющие наименьшее значение метрики расстояния по отношению 


к~референсному товару (в~качестве референсного рассматривается товар, 


который заинтересовал посетителя в~конкретный момент времени). При 


необходимости рекомендации могут быть дополнены предложениями 


информационных единиц, имеющих близкие оценки подобия (метод Item--Item 


CF). Необходимость дополнения рекомендаций, полученных при 


использовании метода User--User CF, будет возникать, как правило, для 


пользователей, не име\-ющих четко выраженных предпочтений. Проведенные 


исследования показывают, что в~целом предложенный подход к~решению 


задачи выработки рекомендаций с~использованием метода User--User CF (на 


базе использования методов кластерного анализа) позволяет формировать 


рекомендации, соответствующие ожиданиям пользователей. Так, например, 


более~70\%~регулярно посещающих ин\-тер\-нет-ма\-га\-зин Thaisoap 


пользователей (по результатам экспериментальных исследований) 


воспользовались предложениями рекомендательного сервиса (совершили 


переходы по предлагаемым сервисом рекомендательным  ссылкам или 


добавляли рекомендованные товары в~корзину).


  


\vspace*{-6pt}





\section{Оценка эффективности метода повышения пертинентности 


информации}





\vspace*{-2pt}





  Одним из наиболее распространенных подходов к~определению 


эффективности методов, применяемых в~рекомендательных системах  


ин\-тер\-нет-ма\-га\-зи\-нов, является оценка средней наполненности 


покупательской корзины до и~после использования рекомендательной системы. 


Существующие исследования дают довольно большой разброс результатов 


измерений. Но сходятся все они на том, что в~любых случаях наблюдается рост 


средней наполненности покупательской корзины, который может 


составлять~12\%--60\%~[9], при этом наблюдается значительная 


дифференциация по типам товаров. Так, для бытовой техники и~электроники 


эффект использования рекомендательных систем, как правило, минимален. Для 


книг и~спортивных товаров, напротив, эффект близок к~максимальному. 


Непосредственно на эффективность влияет такой параметр, как покрытие, 


определяющий, насколько полно имеющийся ассортимент товаров и~услуг 


охватывается при выработке рекомендаций. Косвенное влияние на 


эффективность оказывает качество работы рекомендательной системы. Чем 


более полезные рекомендации вырабатываются, тем больше доверие 


и~лояльность пользователей к~интернет-магазину. Доверие пользователей 


является фактором, который может играть существенную роль в~долгосрочной 


перспективе. 


  


  Для оценки эффективности рекомендательной системы, построенной на базе 


предложенного комбинированного подхода, был проведен сравнительный 


анализ четырех сценариев работы экспериментального образца программного 


комплекса повышения пертинентности информации в~составе  


ин\-тер\-нет-ма\-га\-зи\-на Thaisoap. Первый сценарий не предполагает 


использования рекомендательного сервиса, покупатели используют только 


каталог товаров и~стандартные средства поиска для выбора товаров. Второй 


сценарий предполагает использование рекомендательного сервиса, 


реализующего только метод Item--Item CF. Третий сценарий аналогичен 


второму, только используется метод User--User CF. В~четвертом сценарии 


используется рекомендательный сервис, полностью использующий созданный 


метод повышения пертинентности информации на базе предложенного 


комбинированного подхода (Item--Item CF\;+\;User--User CF). В~каждом из 


рассмотренных сценариев брался временной период, равный одному месяцу. 


При этом каталог товаров оставался неизменным (т.\,е.\ не изменялась 


номенклатура товаров, предлагаемых посетителям магазина). Одним из 


условий проведения сравнительного анализа являлась работа <<с~чистого 


листа>>, т.\,е.\ в~начале каждого месячного отрезка времени не использовались 


данные о пользовательских предпочтениях или рейтингах информационных 


единиц, накопленные за предыдущие периоды времени.


  


  В первом сценарии было выявлено 1860 посетителей, из 


которых~140~чел.\ совершили покупки (конверсия~7,5\%). Средняя 


наполненность корзины при этом составила~1,5~товара. При навигации по 


сайту посетители совершали в~сред\-нем~36~кликов в~рамках одной сессии 


и~проводили на сайте в~среднем порядка~16~мин. Во втором сценарии было 


выявлено~1908~посетителей, из ко\-то\-рых покупки совершили~164~чел.\ 


(конверсия~8,6\%). Средняя наполненность корзины при этом 


составила~2,1~товара. При навигации по сайту посетители совершали 


в~среднем~31~клик в~рамках одной сессии и~проводили на сайте в~среднем 


около~12~мин. В~третьем сценарии было выявлено~1873~посетителя, из 


ко\-то\-рых~145~чел.\ совершили покупки (конверсия~7,7\%). Средняя 


наполненность корзины при этом составила~1,9~товара. При навигации по 


сайту посетители совершали в~среднем~32~клика в~рамках одной сессии 


и~проводили на сайте в~среднем около~14~мин. В~последнем сценарии было 


выявлено~2011~посетителей, из которых~183~чел.\ совершили покупки 


(конверсия~9,1\%). Средняя наполненность корзины при этом 


составила~2,4~товара. При навигации по сайту посетители совершали 


в~среднем~28~ кликов в~рамках одной сессии и~проводили на сайте в~среднем 


около~11~мин. 


  


  Анализ полученных результатов исследования позволяет сделать следующие 


основные выводы.


  \begin{enumerate}[1.]


\item Использование рекомендательной системы в~рассматриваемых 


сценариях позволило повысить среднюю наполненность корзины 


и~сократить время, затрачиваемое посетителями на поиск интересующих их 


товаров. Так, среднее время пребывания посетителей на сайте сокращалось 


в~различных сценариях с~16 до~14--11~мин, при этом в~сценариях 


с~использованием рекомендательной системы увеличивалась средняя 


наполненность корзины, что напрямую влияло на среднюю величину чека 


и~непосредственные доходы магазина. 


\item Метод User--User CF показал себя малоэффективным на 


рассматриваемом интервале времени (один месяц). Основной причиной 


этого стало низкое качество рекомендаций, связанное с~недостаточностью 


данных о пользовательских предпочтениях.


\item Метод Item--Item CF показал себя достаточно эффективным за счет 


успешного решения проблемы <<холодного старта>>. По сравнению 


с~первым сценарием конверсия увеличилась более чем на~1\%, средняя 


наполненность корзины увеличилась с~1,5 до~2,1.


\item Рекомендательная система на базе комбинированного подхода 


позволила получить наибольший прирост рассматриваемых показателей. 


Так, конверсия выросла более чем на~1,5\% по сравнению с~первым 


сценарием, а~средняя наполненность корзины выросла с~1,5 до~2,4 (т.\,е.\ 


более чем на~35\%).


  \end{enumerate}


  


  В целом необходимо отметить, что эффективность рекомендательной 


системы сильно зависит как от типов продаваемых товаров, так и~от 


особенностей используемых алгоритмов и~качества реализации 


рекомендательной системы. В~определенных ситуациях рекомендательная 


система может не давать заметного повышения показателей конверсии 


и~средней наполненности товарной корзины. Тем не менее она остается 


полезной, так как предлагает более удобный и~простой для посетителей способ 


поиска интересующих их товаров, тем самым создавая удобное окружение 


и~повышая лояльность посетителей к~ин\-тер\-нет-ре\-сур\-су, что уже 


в~среднесрочной перспективе может обеспечить возвраты пользователей, 


формирование пула постоянных покупателей, обеспечивающих стабильный 


рост выручки ин\-тер\-нет-ма\-га\-зина.


  


\vspace*{-6pt}





\section{Заключение}





\vspace*{-2pt}





  В данной работе представлено описание метода повышения пертинентности 


информации, основанного на комбинированном использовании методов  


Item--Item CF и~User--User CF с~применением неявных данных. 


  \begin{enumerate}[1.]


\item Для новых пользователей или пользователей с~непрезентабельной 


историей посещений предполагается генерировать рекомендации, базируясь 


на рейтингах информационных единиц (метод Item--Item CF). Таким образом 


решается проблема <<холодного старта>> и~повышается качество 


рекомендаций для малоактивных пользователей, а~также пользователей со 


слабо выраженными предпочтениями. 


\item По мере накопления презентабельных данных о пользовательских 


предпочтениях и~формировании его поведенческого профиля предпочтение 


предлагается отдавать оценкам, полученным с~использованием метода  


User--User CF. При наличии качественных поведенческих профилей 


пользователей данный метод позволяет более точно предсказывать их 


предпочтения. При этом в~случае необходимости рекомендации, полученные 


с использованием метода User--User CF, могут дополняться предложениями 


информационных единиц, полученными на основе их рейтинга 


популярности. 


  \end{enumerate}


  


  Таким образом, предлагаемый метод позволяет более полно использовать 


доступную информацию о пользователях и~информационных единицах, а~также 


нивелировать недостатки обоих методов.


  


  Практическая реализация предложенного метода в~виде экспериментального 


образца программного комплекса повышения пертинентности информации 


экспериментально проверена в~составе ин\-тер\-нет-ма\-га\-зи\-на Thaisoap. 


Анализ журналов активности пользователей в~период экспериментальных 


исследований показал, что более~70\% посетителей магазина используют блоки 


рекомендательной системы для поиска интересующих их товаров. При этом 


процент использования блоков рекомендательной системы посетителями, 


совершившими покупки за указанный период времени, превышает~80\%. 


Сравнительный анализ различных сценариев показывает, что использование 


комбинированного подхода при построении рекомендательной системы 


позволило повысить конверсию более чем на~1,5\%, сократить среднее время 


пребывания посетителей на сайте (с~16 до~11~мин), а~также увеличить 


среднюю наполненность товарной корзины (с~1,5 до~2,4~товара). 


  


  Также необходимо отметить, что предложенный авторами метод 


гарантировано показывает лучшие результаты, чем известные и~применяемые 


на сегодня методы Item--Item CF и~User--User CF. В частности, это можно связать 


с~учетом неявных данных, позволяющих более полно и~быстро выявить 


пользовательские предпочтения, необходимые для формирования 


рекомендаций, что особенно важно сегодня, когда большая часть посетителей 


приходит из поисковой системы и~затем более никогда не возвращается.





\vspace*{-4pt}


  


{\small\frenchspacing


{%\baselineskip=10.8pt


\addcontentsline{toc}{section}{Литература}


\begin{thebibliography}{9}





\vspace*{-2pt}


\bibitem{1-fil}


Почему персонализация контента это еще не веб-персонализация. Блог платформы 


LPgenerator. {\sf 


http://lpgenerator.ru/blog/2016/03/19/pochemu-personalizaciya-kontenta-eto-eshe-ne-veb-personalizaciya}. 


\bibitem{2-fil}


\Au{Филиппов С.\,А., Захаров В.\,Н., Ступников~С.\,А., Ковалев~Д.\,Ю.} Подходы 


к~повышению пертинентности информационного предложения в~медиасервисах на 


основе обработки больших объемов данных~/\!\!/ Сeur Workshop Proceeding: 


Selected Papers of the 17th~Conference (International) on Data Analytics and Management in 


Data Intensive Domains (DAMDID/RCDL 2015), 2015. Vol.~1536. P.~114--118.


\bibitem{3-fil}


\Au{Филиппов С.\,А., Захаров В.\,Н., Ступников~С.\,А., Ковалев~Д.\,Ю.} Организация 


больших объемов данных в~рекомендательных системах поддержки жизнеобеспечения, 


входящих в~состав глобальных платформ электронной коммерции~/\!\!/ Сeur Workshop 


Proceedings: Selected Papers of the 17th~Conference (International) on Data Analytics and 


Management in Data Intensive Domains (DAMDID/RCDL 2015), 2015. Vol.~1536.  


P.~119--124.


\bibitem{4-fil}


\Au{Лексин В.\,А.} Технология персонализации на основе выявления тематических 


профилей пользователей и~ресурсов Интернета.~--- Вычислительный 


центр им. А.\,А.~Дородницина РАН, 2007.  ВКР магистра.


\bibitem{5-fil}


\Au{Su X., Khoshgoftaar T.\,M.} A~survey of collaborative filtering techniques~/\!\!/ Advances 


Artificial Intelligence, 2009. Vol.~2009. Article ID~421425. 19~p.


\bibitem{6-fil}


\Au{Гончаров М.} Системы выработки рекомендаций. {\sf 


http://www.businessdataanalytics.\linebreak ru/RecommendationSystems.htm}.


\bibitem{7-fil}


\Au{Frey B.\,J., Dueck D.} Clustering by passing messages between data points~/\!\!/ Science, 


2007. Vol.~315. Iss.~5814. P.~972--976. doi: 10.1126/science.1136800.


\bibitem{8-fil}


\Au{Coates A., Ng A.\,Y.} Learning feature representations with K-means.~--- Stanford 


University, 2012. {\sf http://www.cs.stanford.edu/$\sim$acoates/papers/coatesng\_nntot2012.pdf}.


\bibitem{9-fil}


\Au{Дьяченко В.} Сервисы рекомендаций: как с~их помощью увеличить продажи 


на~60\%~/\!\!/ Коммерческий директор. {\sf http://www.kom-dir.ru/article/51-servisy-rekomendatsiy}.


       \end{thebibliography}


} }








\vspace*{-3pt}





\hfill{\small\textit{Поступила в~редакцию 14.09.16}}














\vspace*{9pt}





\hrule





%\newpage





\vspace*{2pt}





\hrule





%\vspace*{9pt}











\def\tit{METHOD OF~INCREASING INFORMATION PERTINENCE  


FOR~E-COMMERCE RECOMMENDER SYSTEMS BASED~ON~IMPLICIT 


DATA}





\def\titkol{Method of~increasing information pertinence  


for~e-commerce recommender systems} % based~on~implicit  data}








\def\aut{S.\,A.~Philippov and V.\,N.~Zakharov}


\def\autkol{S.\,A.~Philippov and V.\,N.~Zakharov}











\titel{\tit}{\aut}{\autkol}{\titkol}





\vspace*{-9pt}








\noindent


Institute of Informatics Problems, Federal 


Research Center ``Computer Science and Control'' of the Russian Academy of 


Sciences, 44-2~Vavilov Str., Moscow 119333, Russian Federation








\def\leftfootline{\small{\textbf{\thepage}


\hfill Sistemy i Sredstva Informatiki~---


Systems and Means of Informatics 2016 vol~26 no\,4}


}%


 \def\rightfootline{\small{Sistemy i Sredstva Informatiki~---


 Systems and Means of Informatics   2016   vol~26   no\,4


\hfill \textbf{\thepage}}} 


  


  


\Abste{The paper describes the method of increasing pertinence of information in 


e-commerce recommender systems based on implicit data, i.\,e., due to the processing of 


user activity associated with the decision-making process. The method works 


successfully in situations where information about user activity is absent or little 


informative. Practical application of this method in e-commerce


systems can improve 


their efficiency through targeted supply of goods and services\linebreak\vspace*{-12pt}}





\Abstend{to consumers. The 


main feature of the proposed method is the combined use of Item--Item CF (collaborative


filtering) and  


User--User CF methods taking into account the implicit data collected. The features of the 


proposed method are verified by a prototype software that is installed on the existing 


online store Thaisoap.} 





\KWE{pertinence search; collaborative filtering; e-commerce recommender system; 


implicit data targeting}





  





  


  


\DOI{10.14357\!/\!08696527160401} 





%\vspace*{-24pt}





\Ack


\noindent


This work was supported by the Ministry of Education and Science of the Russian 


Federation, the unique identifier of the project is RFMEFI60414X0139.





\renewcommand{\bibname}{\protect\rmfamily References}


%\renewcommand{\bibname}{\large\protect\rm References}





{\small\frenchspacing


{%\baselineskip=10.8pt


\addcontentsline{toc}{section}{References}


\begin{thebibliography}{9} 


\bibitem{1-fil-1}


Pochemu personalizatsiya kontenta eto eshche ne veb-personalizatsiya [Why 


personalized content is not web personalization]. Blog of the LPgenerator. Available 


at: {\sf  


http://lpgenerator.ru/blog/2016/03/19/pochemu-personalizaciya-kontenta-eto-eshe-ne-veb-personalizaciya/} (accessed August~9, 2106).


\bibitem{2-fil-1}


\Aue{Philippov, S., V. Zakharov, S.~Stupnikov, and D.~Kovalev}. 2015. 


Podkhody k~po\-vy\-she\-niyu pertinentnosti informatsionnogo predlozheniya 


v~mediaservisakh na osno\-ve obrabotki bol'shikh ob"emov dannykh [Approaches 


to improve the pertinence of information in the media services on the basis of big 


data processing]. \textit{Сeur Workshop Proceedings: Selected papers of the 


17th Conference (International) on Data Analytics and Management in Data 


Intensive Domains (DAMDID/RCDL 2015)}. Obninsk, Russia. 1536:114--118. 


Available at:{\sf  


http://ceur-ws.org/Vol-1536/paper17.pdf} (accessed September~29, 2016).


\bibitem{3-fil-1}


\Aue{Philippov, S., V. Zakharov, S.~Stupnikov, and D.~Kovalev}. 2015. 


Organizatsiya bol'shikh ob"emov dannykh v~rekomendatel'nykh sistemakh 


podderzhki zhiz\-ne\-obes\-pe\-che\-niya, vkhodyashchikh v~sostav global'nykh platform 


elektronnoy kommertsii [Organization of big data in the global e-commerce 


platforms]. \textit{Сeur Workshop Proceedings: Selected Papers of the 17th 


Conference (International) on Data Analytics and Management in Data Intensive 


Domains (DAMDID/RCDL 2015)}. Obninsk, Russia. 1536:119--124. Available 


at: {\sf http://ceur-ws.org/Vol-1536/paper18.pdf} (accessed September~29, 2016).


\bibitem{4-fil-1}


\Aue{Leksin, V.\,A.} 2007. Tekhnologiya personalizatsii na osnove vyyavleniya 


tematicheskikh profiley pol'zovateley i~resursov Internet [Personalization 


technology on the basis of the identification of the thematic profiles of users and 


resources online]. Vychislitel'nyy Tsentr im. A.\,A.~Dorodnitsyna RAN. Master's 


Degree Thesis.


\bibitem{5-fil-1}


\Aue{Su, X., and T.\,M.~Khoshgoftaar}. 2009. A~survey of collaborative filtering 


techniques. \textit{Advances Artificial Intelligence} 2009. Article ID~421425. 


19~p.


\bibitem{6-fil-1}


\Aue{Goncharov, M.} 2010. Sistemy vyrabotki rekomendatsiy 


[Recommendations systems]. Available at: 


{\sf http://www.businessdataanalytics.ru/RecommendationSystems.htm} 


(accessed August~12, 2106).


\bibitem{7-fil-1}


\Aue{Frey, B.\,J., and D.~Dueck}. 2007. Clustering by passing messages between 


data points. \textit{Science} 315(5814):972--976. doi: 10.1126/science.1136800.


\bibitem{8-fil-1}


\Aue{Coates, A. and A.\,Y.~Ng}. 2012. \textit{Learning feature representations 


with K-means}. Stanford University. Available at: {\sf 


http://www.cs.stanford.edu/$\sim$acoates/papers/\linebreak coatesng\_nntot2012.pdf} 


(accessed July~14, 2106).


\bibitem{9-fil-1}


\Aue{D'yachenko, V.} 2016. Servisy rekomendatsiy: kak s ikh pomoshch'yu 


uvelichit' prodazhi na~60\% [Recommendation services: How to use them to 


increase the sales by~60\%]. \textit{Kom\-mer\-che\-skiy Director} [Director of 


Marketing]. Available at: 


{\sf http://www.kom-dir.ru/article/51-servisy-rekomendatsiy} (accessed 


August~29, 2106).


\end{thebibliography}


} }








\vspace*{-3pt}





\hfill{\small\textit{Received September 14, 2016}}





  \Contr


  


  \noindent


  \textbf{Philippov Stanislav A.} (b.\ 1980)~--- Candidate of Science (PhD) in 


technology, senior scientist, Institute of Informatics Problems, Federal 


Research Center ``Computer Science and Control'' of the Russian Academy of 


Sciences, 44-2~Vavilov Str., Moscow 119333, Russian Federation; 


\mbox{stanislav@philippov.ru}


  


  \vspace*{3pt}


  


  \noindent


  \textbf{Zakharov Victor N.} (b.\ 1948)~--- Doctor of Science in 


technology, associate professor; Scientific Secretary, Federal Research Center 


``Computer Science and Control'' of the Russian Academy of Sciences, 44-


2~Vavilov Str., Moscow 119333, Russian Federation; 


\mbox{vzakharov@ipiran.ru} 


  


\label{end\stat}








\renewcommand{\bibname}{\protect\rm Литература}


